\section{Annotation Artifacts are Common}
\label{ssec:artifacts}

We conjecture that the framing of the annotation task has a significant effect on the language generation choices that crowd workers make when authoring hypotheses, producing certain patterns in the data.
We call these patterns \textit{annotation artifacts}.


To determine the degree to which such artifacts exist, we train a model to predict the label of a given hypothesis \textit{without seeing the premise}.
Specifically, we use \texttt{fastText} \cite{Joulin:2017}, an off-the-shelf text classifier that models text as a bag of words and bigrams, to predict the entailment label of the hypothesis.\footnote{For MultiNLI, we additionally enabled two hyperparameters: character 4-grams, and filtering words that appeared less than 10 times in the training data.} This classifier is completely oblivious to the premise.


\tabref{results} shows that a significant portion of each test set can be correctly classified without looking at the premise, well beyond the most-frequent-class baseline.\footnote{Experiments with two other text classifiers, a logistic regression classifier with word and character $n$-gram features and a premise-oblivious version of the decomposable attention model \cite{Parikh:2016}, yielded similar results.}

Our finding demonstrates that it is possible to perform well on these datasets without modeling natural language inference.

\begin{table}[t]
\setlength{\tabcolsep}{4.5pt}

\centering
{
\begin{tabular}{lccc}
\toprule
\multirow{2}{*}{\textbf{Model}} & \multirow{2}{*}{\textbf{SNLI}} & \multicolumn{2}{c}{\textbf{MultiNLI}}  \\
 & & \textbf{Matched} & \textbf{Mismatched} \\
\midrule
majority class    & 34.3 & 35.4 & 35.2 \\
%Style                    & 63.0 & 51.4 & 51.2 \\
%NLI (Attention)          & 67.1 & 55.2 & 55.1 \\
% \midrule
\texttt{fastText}        & \bf 67.0 & \bf 53.9 & \bf 52.3 \\
\bottomrule
\end{tabular}
}
\caption{Performance of a premise-oblivious text classifier on NLI. 
The MultiNLI benchmark contains two test sets: matched (in-domain examples) and mismatched (out-of-domain examples). 
A majority baseline is presented for reference.}
\label{tab:results}
\end{table}

